\chapter{Greeks and the Hedging P\&L}\label{ch:dv:greeks-and-the-hedging-pnl}

An options desk sells one-month at-the-money \omterm{def:dv:greeks-and-the-hedging-pnl:cashgamma}{straddles} on a share at 20
volatility and hedges their delta every evening. For the next four weeks the
share moves more than usual; its realised volatility over the month is 25. On
the first morning the desk's risk report printed two numbers, a \omterm{def:dv:greeks-and-the-hedging-pnl:greeks}{vega} of
USD~23\,000 per volatility point and a \omterm{def:dv:greeks-and-the-hedging-pnl:greeks}{theta} of USD~7\,600 a day in its favour,
and one line of arithmetic: if the share realises five points more than was
sold, the book loses about five times its \omterm{def:dv:greeks-and-the-hedging-pnl:greeks}{vega}, USD~115\,000. At the end of the
month the loss is close to that figure, and the desk can say, day by day, where
it came from: on each day the gamma of the position paid out more than the
\omterm{def:dv:greeks-and-the-hedging-pnl:greeks}{theta} brought in. This chapter makes that bookkeeping exact. It names the
sensitivities, proves the identity that links gamma to \omterm{def:dv:greeks-and-the-hedging-pnl:greeks}{theta}, turns it into a
formula for the P\&L of a hedged option, and measures what daily rather than
continuous hedging adds to it.

\section{The Greeks and their shapes}

The delta and the gamma of an option are defined in One Quant Book 1, chapter 26:
the first and second derivatives of its value with respect to the underlying.
Delta hedging holds minus the delta in the underlying.

\begin{definition}[Greeks]\label{def:dv:greeks-and-the-hedging-pnl:greeks}
The \emph{Greeks}\index{Greeks} of a position are the partial derivatives of its
value $V(t,S,\sigma,r)$ with respect to the inputs of the pricing model: delta
$\Delta=\partial_SV$ and gamma $\Gamma=\partial_{SS}V$; the \emph{vega}\index{vega}
$\mathcal V=\partial_\sigma V$; the \emph{theta}\index{theta} $\Theta=\partial_tV$, the
change of value with calendar time, all else fixed; the \emph{rho}\index{rho}
$\mathrm{Rho}=\partial_rV$; and the second-order cross-sensitivities
\emph{vanna}\index{vanna} $\mathrm{Vanna}=\partial_{S\sigma}V$, the change of delta
with volatility, and \emph{volga}\index{volga} $\mathrm{Volga}=\partial_{\sigma\sigma}V$,
the change of vega with volatility.
\end{definition}

For a call in the \omterm{def:dv:black-scholes-three-ways:model}{Black--Scholes model}, with $\varphi$ the normal density,
\[
\Gamma=\frac{e^{-qT}\varphi(d_1)}{S\sigma\sqrt T},\quad
\mathcal V=Se^{-qT}\varphi(d_1)\sqrt T,\quad
\mathrm{Vanna}=-\frac{e^{-qT}\varphi(d_1)\,d_2}{\sigma},\quad
\mathrm{Volga}=\mathcal V\,\frac{d_1d_2}{\sigma},
\]
and the put has the same gamma and \omterm{def:dv:greeks-and-the-hedging-pnl:greeks}{vega}. Desks report \omterm{def:dv:greeks-and-the-hedging-pnl:greeks}{Greeks} in money: \omterm{def:dv:greeks-and-the-hedging-pnl:greeks}{vega} per
volatility point (0.01), \omterm{def:dv:greeks-and-the-hedging-pnl:greeks}{theta} per calendar day, \omterm{def:dv:greeks-and-the-hedging-pnl:greeks}{rho} per basis point, and gamma
through the next quantity.

\begin{definition}[Cash gamma and straddle]\label{def:dv:greeks-and-the-hedging-pnl:cashgamma}
The \emph{cash gamma}\index{cash gamma} of a position is $\tfrac12\Gamma S^2$: the P\&L
of the delta-hedged position per unit of squared return, since
$\tfrac12\Gamma\,dS^2=\tfrac12\Gamma S^2(dS/S)^2$. Risk reports also give
$\Gamma S^2/100$, the change of the cash delta $\Delta S$ for a 1\% move. A
\emph{straddle}\index{straddle} is a call and a put with the same strike and expiry:
at the money its delta is close to zero and its gamma and \omterm{def:dv:greeks-and-the-hedging-pnl:greeks}{vega} are twice the
call's.
\end{definition}

\begin{omfigure}
\begin{tikzpicture}
\begin{groupplot}[group style={group size=2 by 1,horizontal sep=1.6cm},
  width=6.6cm,height=4.8cm,xlabel={spot},
  tick label style={font=\small},label style={font=\small},
  xmin=70,xmax=130,grid=major,grid style={gray!25},table/col sep=comma]
\nextgroupplot[ylabel={gamma},ymin=0,ymax=0.095,
  legend columns=3,legend style={font=\footnotesize,at={(1.15,-0.32)},anchor=north,draw=none,fill=none,/tikz/every even column/.append style={column sep=8pt}}]
\addplot[omThm,very thick] table[x=spot,y=g1]{figdata/derivatives/04-greeks-and-the-hedging-pnl/shapes.csv};
\addplot[omDef,thick] table[x=spot,y=g025]{figdata/derivatives/04-greeks-and-the-hedging-pnl/shapes.csv};
\addplot[omProp,thick] table[x=spot,y=g005]{figdata/derivatives/04-greeks-and-the-hedging-pnl/shapes.csv};
\legend{1 year,3 months,2.5 weeks}
\nextgroupplot[ylabel={vega per point},ymin=0,ymax=0.42]
\addplot[omThm,very thick] table[x=spot,y=v1]{figdata/derivatives/04-greeks-and-the-hedging-pnl/shapes.csv};
\addplot[omDef,thick] table[x=spot,y=v025]{figdata/derivatives/04-greeks-and-the-hedging-pnl/shapes.csv};
\addplot[omProp,thick] table[x=spot,y=v005]{figdata/derivatives/04-greeks-and-the-hedging-pnl/shapes.csv};
\end{groupplot}
\end{tikzpicture}

\omcaption{Gamma (left) and \omterm{def:dv:greeks-and-the-hedging-pnl:greeks}{vega} per volatility point (right) of a call struck at
100 ($\sigma=20\%$, $r=0$) at three expiries. Short options concentrate gamma
near the strike; long options carry the \omterm{def:dv:greeks-and-the-hedging-pnl:greeks}{vega}. Both are positive for every long
option. Data: the chapter's code.}\label{fig:dv:greeks-and-the-hedging-pnl:shapes}
\end{omfigure}

\section{The gamma--theta identity and the hedging P\&L}

\begin{proposition}[Gamma--theta identity]\label{prop:dv:greeks-and-the-hedging-pnl:gt}
The value of any European claim in the \omterm{def:dv:black-scholes-three-ways:model}{Black--Scholes model} satisfies
\[
\Theta+\tfrac12\sigma^2S^2\Gamma=r(V-S\Delta)+qS\Delta .
\]
With zero rates and dividends, $\Theta=-\tfrac12\sigma^2S^2\Gamma$: a long option
loses in time exactly what its gamma earns when the share moves by one standard
deviation.
\end{proposition}
\begin{proof}
Rearrange the \omterm{def:dv:black-scholes-three-ways:pde}{Black--Scholes equation} of chapter 3.
\end{proof}

Now hold the option priced and hedged at a volatility $\sigma$, while the share
moves as it will. Over a short interval the position (option, $-\Delta$ shares,
cash financing both) changes by a Taylor expansion in which the delta term is
hedged away:
\[
d\Pi=\Theta\,dt+\tfrac12\Gamma\,dS^2-r(V-S\Delta)\,dt-qS\Delta\,dt
=\tfrac12\Gamma S^2\Bigl(\bigl(\tfrac{dS}{S}\bigr)^2-\sigma^2dt\Bigr),
\]
by the identity. This is the whole theory of a volatility book in one line.

\begin{definition}[Hedging P\&L]\label{def:dv:greeks-and-the-hedging-pnl:hpnl}
The \emph{hedging P\&L}\index{hedging P\&L} of a delta-hedged option is the change
of value of the option together with its hedge and its financing. For a long
option hedged at volatility $\sigma$ it is, to second order over each interval,
the \omterm{def:dv:greeks-and-the-hedging-pnl:cashgamma}{cash gamma} times the difference between the squared return and the variance
the hedge assumed:
\[
\pnl_{[0,T]}=\int_0^Te^{r(T-t)}\tfrac12\Gamma_tS_t^2\bigl(\sigma_{\mathrm{real},t}^2-\sigma^2\bigr)dt ,
\]
where $\sigma^2_{\mathrm{real},t}dt=(dS_t/S_t)^2$ is the instantaneous realised
variance.
\end{definition}

The formula says three things. The P\&L is paid at the rate of the \omterm{def:dv:greeks-and-the-hedging-pnl:cashgamma}{cash gamma},
so moves on days when gamma is large count more than moves on days when it is
small. It is quadratic in returns and blind to their sign. And it involves
realised variance, not realised volatility, a fact that chapter 14 turns into a
product.

\begin{proposition}[Expected hedging P\&L]\label{prop:dv:greeks-and-the-hedging-pnl:expected}
If the share follows the \omterm{def:dv:black-scholes-three-ways:model}{Black--Scholes model} with volatility $\sigma_{\mathrm{real}}$
and the option is sold at $\sigma$ and hedged continuously at $\sigma$, the
expected \omterm{def:dv:greeks-and-the-hedging-pnl:hpnl}{hedging P\&L} of the long option (with the drift at $r$) is
$V(\sigma_{\mathrm{real}})-V(\sigma)$ in present value.
\end{proposition}
\begin{proof}
The present value of the P\&L is $\E\bigl[\int_0^Te^{-rt}\tfrac12\Gamma_t(\sigma)S_t^2
(\sigma_{\mathrm{real}}^2-\sigma^2)dt\bigr]$. The function $u=V(\sigma_{\mathrm{real}})-
V(\sigma)$ solves the pricing equation at $\sigma_{\mathrm{real}}$ with the source
term $\tfrac12(\sigma_{\mathrm{real}}^2-\sigma^2)S^2\Gamma(\sigma)$ and zero terminal
value; the Feynman--Kac formula (One Quant Book 4, chapter 4) gives the claim.
\end{proof}

\begin{example}[The desk's month]\label{ex:dv:greeks-and-the-hedging-pnl:month}
The one-month at-the-money \omterm{def:dv:greeks-and-the-hedging-pnl:cashgamma}{straddle} on a share at 100 is worth 4.6059 at 20
volatility and 5.7570 at 25, with zero rates. A desk short 1\,000 \omterm{def:dv:greeks-and-the-hedging-pnl:cashgamma}{straddles}
(multiplier 100) expects to lose $(5.7570-4.6059)\times100\,000$, about
USD~115\,100, when realised volatility comes in at 25. Its \omterm{def:dv:greeks-and-the-hedging-pnl:greeks}{vega}, 0.2302 per point
per \omterm{def:dv:greeks-and-the-hedging-pnl:cashgamma}{straddle}, times five points gives nearly the same number: for small changes
the expected loss is \omterm{def:dv:greeks-and-the-hedging-pnl:greeks}{vega} times the volatility gap.
\end{example}

\begin{omfigure}
\begin{tikzpicture}
\begin{axis}[width=10.5cm,height=5cm,
  xlabel={trading days},ylabel={P\&L per straddle},
  tick label style={font=\small},label style={font=\small},
  xmin=0,xmax=21,ymin=-2,ymax=0.4,grid=major,grid style={gray!25},
  legend columns=2,legend style={font=\footnotesize,at={(0.5,-0.3)},anchor=north,draw=none,fill=none,/tikz/every even column/.append style={column sep=8pt}},
  table/col sep=comma]
\addplot[omDef,very thick] table[x=day,y=pnl]{figdata/derivatives/04-greeks-and-the-hedging-pnl/gamma_path.csv};
\addplot[omThm,thick,dashed] table[x=day,y=pred]{figdata/derivatives/04-greeks-and-the-hedging-pnl/gamma_path.csv};
\legend{hedged short straddle,$-\sum\frac12\Gamma S^2\bigl(r_k^2-\sigma^2\Delta t\bigr)$}
\end{axis}
\end{tikzpicture}

\omcaption{One month of a short \omterm{def:dv:greeks-and-the-hedging-pnl:cashgamma}{straddle} sold at 20 volatility, hedged four times
a day, while the share realises 25. The dashed line adds the gamma--theta
terms interval by interval; it never strays more than 0.09 from the book's P\&L
and ends 0.04 from it.
The P\&L falls in steps on the days with large moves and drifts up on quiet days,
when \omterm{def:dv:greeks-and-the-hedging-pnl:greeks}{theta} wins. Data: the tutorial.}\label{fig:dv:greeks-and-the-hedging-pnl:path}
\end{omfigure}

\section{Discrete hedging error}

The identity is exact only in the limit of continuous hedging. A desk that hedges
once a day holds a stale delta between rebalancings, and the squared return over
a day is a noisy estimate of the day's variance.

\begin{definition}[Discrete hedging error]\label{def:dv:greeks-and-the-hedging-pnl:dhe}
The \emph{discrete hedging error}\index{discrete hedging error} is the difference
between the terminal value of an option hedged at $N$ dates and that of the same
option hedged continuously, when realised volatility equals the hedging volatility.
Its mean is close to zero and its standard deviation falls like $1/\sqrt N$.
\end{definition}

\begin{proposition}[Size of the error]\label{prop:dv:greeks-and-the-hedging-pnl:dk}
For an option near the money hedged $N$ times at its implied volatility $\sigma$,
which the share realises, the standard deviation of the hedging error is
approximately $\sqrt{\pi/4}\,\mathcal V\sigma/\sqrt N$.
\end{proposition}
\begin{proof}[Partial proof]
Over a step the error is $\tfrac12\Gamma S^2(x^2-\sigma^2\Delta t)$ with $x$ the
step's return; its variance is $\tfrac12\Gamma^2S^4\sigma^4\Delta t^2$ for normal
returns. Summing the independent steps, $\Delta t=T/N$, and averaging
$\Gamma^2S^4$ over the paths gives the $1/\sqrt N$ rate and the \omterm{def:dv:greeks-and-the-hedging-pnl:greeks}{vega} scaling; the
constant $\sqrt{\pi/4}$ comes from that average for an at-the-money option
(Derman and Kamal).
\end{proof}

For the one-month \omterm{def:dv:greeks-and-the-hedging-pnl:cashgamma}{straddle} hedged daily, the formula gives 0.890 and a simulation
of 20\,000 paths gives 0.865; hedging eight times a day brings the simulated
figure to 0.308 (\cref{fig:dv:greeks-and-the-hedging-pnl:discrete}). The noise is
as large as the edge: a desk that sells volatility two points above what the
share will realise earns, on average, twice the \omterm{def:dv:greeks-and-the-hedging-pnl:greeks}{vega} per \omterm{def:dv:greeks-and-the-hedging-pnl:cashgamma}{straddle}, about 0.46, and
with daily hedging has a standard deviation of 0.87 around it. Selling volatility
is a statistical business, run on many positions.

\begin{omfigure}
\begin{tikzpicture}
\begin{axis}[width=10.5cm,height=4.8cm,
  xlabel={P\&L at expiry per straddle},ylabel={density},
  tick label style={font=\small},label style={font=\small},
  xmin=-4,xmax=4,ymin=0,ymax=1.7,grid=major,grid style={gray!25},
  legend columns=2,legend style={font=\footnotesize,at={(0.5,-0.3)},anchor=north,draw=none,fill=none,/tikz/every even column/.append style={column sep=8pt}},
  table/col sep=comma]
\addplot[omDef,thick,const plot mark mid] table[x=x,y=daily]{figdata/derivatives/04-greeks-and-the-hedging-pnl/discrete.csv};
\addplot[omThm,thick,const plot mark mid] table[x=x,y=eight]{figdata/derivatives/04-greeks-and-the-hedging-pnl/discrete.csv};
\legend{hedged daily (21 times),hedged 8 times a day (168)}
\end{axis}
\end{tikzpicture}

\omcaption{Distribution of the P\&L of a short one-month \omterm{def:dv:greeks-and-the-hedging-pnl:cashgamma}{straddle} sold at 20
volatility when the share realises exactly 20, over 20\,000 simulated paths. Both
distributions are centred on zero; their standard deviations, 0.85 and 0.31, fall
like one over the square root of the number of hedges. Data: the
tutorial.}\label{fig:dv:greeks-and-the-hedging-pnl:discrete}
\end{omfigure}

\section{Break-even volatility}

\begin{definition}[Break-even volatility]\label{def:dv:greeks-and-the-hedging-pnl:be}
The \emph{break-even volatility}\index{break-even volatility} of a hedged option is
the realised volatility at which, over a period, its gamma P\&L exactly pays its
\omterm{def:dv:greeks-and-the-hedging-pnl:greeks}{theta} and financing. The break-even move is the corresponding daily move of the
underlying: $\tfrac12\Gamma\,\delta S^2=-\Theta\,\delta t$.
\end{definition}

With zero rates the gamma--theta identity makes the \omterm{def:dv:greeks-and-the-hedging-pnl:be}{break-even volatility} equal to
the implied volatility at which the option is marked. For the \omterm{def:dv:greeks-and-the-hedging-pnl:cashgamma}{straddle} of the
example, \omterm{def:dv:greeks-and-the-hedging-pnl:greeks}{theta} is $-0.0757$ per calendar day and gamma 0.1381: the break-even move
is $\sqrt{2\times0.0757/0.1381}=1.047$, a day's standard deviation at 20
volatility over 365 days. Desks that count \omterm{def:dv:greeks-and-the-hedging-pnl:greeks}{theta} over calendar days but see moves
only on trading days need the move on trading days to pay the weekend's \omterm{def:dv:greeks-and-the-hedging-pnl:greeks}{theta} as
well: over 252 trading days the break-even move is 1.26. The choice of clock is
a modelling decision, taken up in chapter 8.

\section{Which volatility to hedge at}

A desk that is sure the share will realise 25 still chooses the volatility in its
delta. Two choices are natural.

\begin{proposition}[Hedging at implied against hedging at realised]\label{prop:dv:greeks-and-the-hedging-pnl:which}
Suppose the option is bought at implied $\sigma$ and the share realises a constant
$\sigma_{\mathrm{real}}>\sigma$. Hedging continuously with the delta at
$\sigma_{\mathrm{real}}$ locks in the profit $V(\sigma_{\mathrm{real}})-V(\sigma)$ at
inception, but the mark-to-market at implied fluctuates along the way. Hedging with
the delta at $\sigma$ makes each day's P\&L a known function of that day's return
$x$, namely $\tfrac12\Gamma S^2(x^2-\sigma^2\delta t)$, but the total depends on
the path through the gamma weighting; its mean is the same.
\end{proposition}
\begin{proof}
With the delta at $\sigma_{\mathrm{real}}$, the hedged option is a \omterm{def:dv:no-arbitrage-and-the-fundamental-theorems:claim}{replicating
portfolio} for the true model: its terminal value is fixed and equals its value at
$\sigma_{\mathrm{real}}$ at inception. With the delta at $\sigma$, apply
\cref{def:dv:greeks-and-the-hedging-pnl:hpnl} and
\cref{prop:dv:greeks-and-the-hedging-pnl:expected}.
\end{proof}

In practice realised volatility is not known in advance, so the second choice is
the default: the hedge is at implied, the book's P\&L is explained by gamma,
\omterm{def:dv:greeks-and-the-hedging-pnl:greeks}{theta} and \omterm{def:dv:greeks-and-the-hedging-pnl:greeks}{vega} every day, and the realised-volatility view is expressed in the size
of the position. \Cref{fig:dv:greeks-and-the-hedging-pnl:which} shows both choices
on the desk's month, hedged four times a day: same mean, a standard deviation of
0.72 at implied and 0.54 at realised, the latter being only discrete hedging noise.

\begin{omfigure}
\begin{tikzpicture}
\begin{axis}[width=10.5cm,height=4.8cm,
  xlabel={P\&L at expiry per straddle},ylabel={density},
  tick label style={font=\small},label style={font=\small},
  xmin=-5,xmax=3,ymin=0,ymax=1.0,grid=major,grid style={gray!25},
  legend columns=2,legend style={font=\footnotesize,at={(0.5,-0.3)},anchor=north,draw=none,fill=none,/tikz/every even column/.append style={column sep=8pt}},
  table/col sep=comma]
\addplot[omDef,thick,const plot mark mid] table[x=x,y=implied]{figdata/derivatives/04-greeks-and-the-hedging-pnl/which_vol.csv};
\addplot[omThm,thick,const plot mark mid] table[x=x,y=realised]{figdata/derivatives/04-greeks-and-the-hedging-pnl/which_vol.csv};
\draw[gray,dashed] (axis cs:-1.151,0) -- (axis cs:-1.151,1.0);
\legend{hedged at implied 20,hedged at realised 25}
\end{axis}
\end{tikzpicture}

\omcaption{A short \omterm{def:dv:greeks-and-the-hedging-pnl:cashgamma}{straddle} sold at 20 while the share realises 25, hedged four
times a day with the delta at 20 or at 25. Both centre on the expected loss
$-1.151$ (dashed). Hedging at the true volatility removes the path dependence and
leaves only the discrete-hedging noise. Data: the
tutorial.}\label{fig:dv:greeks-and-the-hedging-pnl:which}
\end{omfigure}

\begin{definition}[Bump-and-reprice]\label{def:dv:greeks-and-the-hedging-pnl:bump}
\emph{Bump-and-reprice}\index{bump-and-reprice} computes a Greek by repricing the
position with one input shifted and taking a finite difference: delta as
$(V(S+h)-V(S-h))/2h$, gamma as $(V(S+h)-2V(S)+V(S-h))/h^2$, \omterm{def:dv:greeks-and-the-hedging-pnl:greeks}{vega} with a shift of
one volatility point. It works for any pricer, at the cost of two or three
valuations per Greek.
\end{definition}

The step is a trade-off (One Quant Book 4, chapter 25): too large and the
difference measures a chord, not a slope; too small and catastrophic cancellation
leaves only rounding noise, especially for gamma. Desks use a relative spot step of
about 1\% for risk reports, matching the moves they care about, and a smaller one
when a smooth derivative is needed. For Monte Carlo pricers the repricings must
use the same random numbers, or the noise of two independent estimates swamps the
difference (chapter 23).

\section{Tutorial: a month of hedging}\label{tut:dv:greeks-and-the-hedging-pnl:tut}

\begin{tutorial}
\textbf{Goal.} Simulate the desk's month in full: sell \omterm{def:dv:greeks-and-the-hedging-pnl:cashgamma}{straddles}, hedge on a
schedule, compare the P\&L with the gamma--theta prediction, and measure the
discrete-hedging noise and the effect of the hedging volatility. \textbf{End state:}
\cref{fig:dv:greeks-and-the-hedging-pnl:path,fig:dv:greeks-and-the-hedging-pnl:discrete,fig:dv:greeks-and-the-hedging-pnl:which}
and the numbers of the weekend problem.
\begin{tutsteps}
\item \textbf{The simulator}, vectorised over 20\,000 paths: sell at implied, hedge
at a chosen volatility, finance at $r$, settle the \omterm{def:dv:greeks-and-the-hedging-pnl:cashgamma}{straddle} at expiry.
\omcode{code/derivatives/04-greeks-and-the-hedging-pnl/python/dv_greeks.py}{46}{69}{Delta-hedging a short straddle on many paths at once.}\label{lst:dv:greeks-and-the-hedging-pnl:sim}
\item \textbf{\omterm{def:dv:greeks-and-the-hedging-pnl:greeks}{Greeks} in desk units} by \omterm{def:dv:greeks-and-the-hedging-pnl:bump}{bump-and-reprice}, for any pricer:
\omcode{code/firm/greeks/firm_greeks.py}{16}{30}{Central bump-and-reprice Greeks, reported per point, per day and per basis point.}\label{lst:dv:greeks-and-the-hedging-pnl:bump}
\item \textbf{Run} \texttt{dv\_greeks.problem()} and \texttt{fig\_greeks.py}; compare
the simulated standard deviations with \cref{prop:dv:greeks-and-the-hedging-pnl:dk}.
\end{tutsteps}
\textbf{What to change next.} Add a bid--ask cost on every hedge trade and find the
hedging frequency that minimises cost plus noise; then let volatility be 25 in the
first half of the month and 15 in the second, and compare the P\&L of a \omterm{def:dv:greeks-and-the-hedging-pnl:cashgamma}{straddle}
struck at 100 on paths that end at 100 and at 110.
\end{tutorial}

\section{Build: the Greek calculator}\label{bld:dv:greeks-and-the-hedging-pnl:greeks}

\begin{build}
\bfield{Purpose} Every position of the miniature firm reports the same \omterm{def:dv:greeks-and-the-hedging-pnl:greeks}{Greeks} in
the same units, whatever pricer values it; the next day's P\&L is predicted from
them, so that the unexplained part can be monitored (chapter 25).
\bfield{Interface} \texttt{bump\_greeks(pricer, params, h\_spot, h\_vol, h\_rate,
days)} with \texttt{pricer(**params)}; \texttt{desk\_units(greeks, spot)};
\texttt{predict\_pnl(greeks, d\_spot, d\_vol\_pts, days)};
\texttt{hedging\_pnl(cash\_gamma, ret, vol, dt)}; \texttt{breakeven\_vol}.
\bfield{Rules} Central differences for delta, gamma, \omterm{def:dv:greeks-and-the-hedging-pnl:greeks}{vega} and \omterm{def:dv:greeks-and-the-hedging-pnl:greeks}{rho}; \omterm{def:dv:greeks-and-the-hedging-pnl:greeks}{theta} by rolling the
valuation date one calendar day; units: shares, $\Gamma S^2/100$, per volatility
point, per day, per basis point.
\bfield{Acceptance tests} \texttt{code/firm/greeks/tests/}: \omterm{def:dv:greeks-and-the-hedging-pnl:bump}{bump-and-reprice} equals
the analytic \omterm{def:dv:greeks-and-the-hedging-pnl:greeks}{Greeks} of \texttt{firm.bs} in desk units; the one-day prediction is
within one cent of a full revaluation for a 1\% move; the \omterm{def:dv:greeks-and-the-hedging-pnl:be}{break-even volatility}
equals the implied volatility at zero rates.
\bfield{Stretch} The same \omterm{def:dv:greeks-and-the-hedging-pnl:greeks}{Greeks} for a portfolio priced by the Monte Carlo engine of
chapter 23, with common random numbers; the pricing library of chapter 28 serves
them for any instrument through one call.
\end{build}

\begin{omsources}
\item P.~P.~Boyle and D.~Emanuel, ``Discretely adjusted option hedges'',
\emph{Journal of Financial Economics} 8 (1980) 259--282.
\item E.~Derman and M.~Kamal, ``When you cannot hedge continuously: the
corrections of Black--Scholes'', \emph{Risk} 12 (1999) 82--85.
\item R.~Ahmad and P.~Wilmott, ``Which free lunch would you like today, sir?
Delta hedging, volatility arbitrage and optimal portfolios'', \emph{Wilmott}
(2005) 64--79.
\item H.~E.~Leland, ``Option pricing and replication with transactions costs'',
\emph{Journal of Finance} 40 (1985) 1283--1301.
\end{omsources}

\section{Exercises}

\begin{exercise}[$\star$]\label{exo:dv:greeks-and-the-hedging-pnl:1}
Give the delta, gamma, \omterm{def:dv:greeks-and-the-hedging-pnl:greeks}{vega} per point and \omterm{def:dv:greeks-and-the-hedging-pnl:greeks}{theta} per day of the one-month
at-the-money \omterm{def:dv:greeks-and-the-hedging-pnl:cashgamma}{straddle} on a share at 100 at 20 volatility, zero rates.
\end{exercise}

\begin{exercise}[$\star$]\label{exo:dv:greeks-and-the-hedging-pnl:2}
Check the gamma--theta identity on the numbers of \cref{exo:dv:greeks-and-the-hedging-pnl:1}.
\end{exercise}

\begin{exercise}[$\star$]\label{exo:dv:greeks-and-the-hedging-pnl:3}
A hedged long option has a \omterm{def:dv:greeks-and-the-hedging-pnl:cashgamma}{cash gamma} of 5\,000 and is marked at 20 volatility.
The share moves 2\% today. Give the day's gamma P\&L, the \omterm{def:dv:greeks-and-the-hedging-pnl:greeks}{theta}, and the net (zero
rates, one calendar day).
\end{exercise}

\begin{exercise}[$\star\star$]\label{exo:dv:greeks-and-the-hedging-pnl:4}
Why is the \omterm{def:dv:greeks-and-the-hedging-pnl:hpnl}{hedging P\&L} blind to the sign of the moves? Give a path on which a
long \omterm{def:dv:greeks-and-the-hedging-pnl:cashgamma}{straddle} loses money although the share ends far from the strike.
\end{exercise}

\begin{exercise}[$\star\star$]\label{exo:dv:greeks-and-the-hedging-pnl:5}
By how much must the number of hedges grow to halve the discrete-hedging error?
What does it cost if each hedge trade pays half a spread?
\end{exercise}

\begin{exercise}[$\star\star$]\label{exo:dv:greeks-and-the-hedging-pnl:6}
Using \cref{prop:dv:greeks-and-the-hedging-pnl:dk}, give the standard deviation of
the daily-hedging error of the desk's \omterm{def:dv:greeks-and-the-hedging-pnl:cashgamma}{straddle}, and compare it with the simulated
0.865.
\end{exercise}

\begin{exercise}[$\star\star\star$]\label{exo:dv:greeks-and-the-hedging-pnl:7}
\emph{Coding.} With \texttt{hedge\_short\_straddle}, measure the mean and standard
deviation of the P\&L when the \omterm{def:dv:greeks-and-the-hedging-pnl:cashgamma}{straddle} is sold at 20, the share realises 25 and the
hedge is at 25, hedged daily, on seed 1.
\end{exercise}

\begin{exercise}[$\star\star\star$]\label{exo:dv:greeks-and-the-hedging-pnl:8}
\emph{Find the flaw.} ``Our short \omterm{def:dv:greeks-and-the-hedging-pnl:cashgamma}{straddles} made money on 70\% of the days this month
and lost overall. \omterm{def:dv:greeks-and-the-hedging-pnl:greeks}{Theta} was on our side; the loss must be a pricing error.''
\end{exercise}

\section{Problem: The Short-Straddle Month}

\begin{problem}[{Weekend problem --- selling volatility that turned out too cheap}]\label{pb:dv:greeks-and-the-hedging-pnl:1}
A desk sells 1\,000 one-month at-the-money \omterm{def:dv:greeks-and-the-hedging-pnl:cashgamma}{straddles} (multiplier 100) on a share at
100, at 20 implied volatility, zero rates and no dividend, and hedges daily. Over the
month the share realises 25.

\medskip\noindent\textbf{Part I --- The position.}
\begin{enumerate}
\item Give the premium received, per \omterm{def:dv:greeks-and-the-hedging-pnl:cashgamma}{straddle} and in dollars.
\item Give the position's delta in shares and the initial hedge.
\item Give its \omterm{def:dv:greeks-and-the-hedging-pnl:cashgamma}{cash gamma} and its gamma per 1\% move, per \omterm{def:dv:greeks-and-the-hedging-pnl:cashgamma}{straddle}.
\item Give its \omterm{def:dv:greeks-and-the-hedging-pnl:greeks}{vega} in dollars per volatility point and its \omterm{def:dv:greeks-and-the-hedging-pnl:greeks}{theta} per day.
\item Give the break-even daily move over 365 and over 252 days.
\end{enumerate}

\medskip\noindent\textbf{Part II --- The expected outcome.}
\begin{enumerate}[resume]
\item Give the \omterm{def:dv:greeks-and-the-hedging-pnl:cashgamma}{straddle}'s value at 25 volatility.
\item Give the expected loss per \omterm{def:dv:greeks-and-the-hedging-pnl:cashgamma}{straddle} and in dollars.
\item Compare it with five times the \omterm{def:dv:greeks-and-the-hedging-pnl:greeks}{vega}.
\item Why do the two agree so closely?
\item Would the answer change if the desk hedged continuously?
\end{enumerate}

\medskip\noindent\textbf{Part III --- The simulated month.}
\begin{enumerate}[resume]
\item Give the simulated mean and standard deviation of the P\&L per \omterm{def:dv:greeks-and-the-hedging-pnl:cashgamma}{straddle} with daily hedging.
\item Give the standard deviation with eight hedges a day.
\item Give the probability of a loss with daily hedging.
\item Give the mean and standard deviation when the desk hedges at 25 instead of 20.
\item With realised volatility equal to 20, what are the mean and standard deviation?
\end{enumerate}

\medskip\noindent\textbf{Part IV --- Judgement.}
\begin{enumerate}[resume]
\item What would the desk have needed to believe to sell at 20?
\item How would you size the position so that one bad month does not end the desk?
\item Which Greek of the report should have warned the desk?
\item State the \emph{named result}: the expected loss of the position from the
gamma--theta identity and its standard deviation under daily hedging, in dollars.
\item In one sentence: what does a delta-hedged option pay?
\end{enumerate}
\end{problem}

\section{Interview questions}

\begin{interviewq}[$\star$ \iqroles{trader}]\label{iq:dv:greeks-and-the-hedging-pnl:1}
You are long a delta-hedged at-the-money option. The underlying does not move for
a week. What happened to your P\&L, and why?
\end{interviewq}

\begin{interviewq}[$\star$ \iqroles{trader, researcher}]\label{iq:dv:greeks-and-the-hedging-pnl:2}
State the gamma--theta relation and explain it without formulas.
\end{interviewq}

\begin{interviewq}[$\star\star$ \iqroles{researcher}]\label{iq:dv:greeks-and-the-hedging-pnl:3}
Derive the P\&L of a delta-hedged option over one day in terms of its gamma and
the realised and implied volatilities.
\end{interviewq}

\begin{interviewq}[$\star\star$ \iqroles{trader}]\label{iq:dv:greeks-and-the-hedging-pnl:4}
You are sure realised volatility will be higher than implied. Should you delta-hedge
at implied or at your forecast? What changes?
\end{interviewq}

\begin{interviewq}[$\star\star$ \iqroles{developer, risk}]\label{iq:dv:greeks-and-the-hedging-pnl:5}
Your \omterm{def:dv:greeks-and-the-hedging-pnl:bump}{bump-and-reprice} gamma is noisy for a Monte Carlo pricer and wrong for a
digital near expiry. Why, and what do you do?
\end{interviewq}

\begin{interviewq}[$\star\star\star$ \iqroles{researcher, trader}]\label{iq:dv:greeks-and-the-hedging-pnl:6}
Two books have the same \omterm{def:dv:greeks-and-the-hedging-pnl:greeks}{vega} and the same realised-minus-implied volatility over a
month. Why can their P\&Ls differ, and by how much?
\end{interviewq}
